\documentclass[11pt]{article}
\usepackage{amsmath,amssymb,amsthm}
\usepackage[margin=1in]{geometry}
\usepackage{hyperref}

\newtheorem{definition}{Definition}[section]
\newtheorem{theorem}[definition]{Theorem}
\newtheorem{lemma}[definition]{Lemma}
\newtheorem{corollary}[definition]{Corollary}
\newtheorem{proposition}[definition]{Proposition}
\newtheorem{remark}[definition]{Remark}

\title{The Logarithm-Vector Isomorphism}
\author{Formalization of ideas from ``Everything Is Logarithms'' by Alex Kritchevsky}
\date{\today}

\begin{document}
\maketitle

\begin{abstract}
We establish a precise correspondence between the baseless logarithm and geometric vectors in differential geometry. We show that logarithms behave as vectors under a natural group action, that vectors in differential geometry can be interpreted as logarithms of translation operators, and that the change-of-base formula for logarithms is isomorphic to coordinate transformations for vectors.
\end{abstract}

\tableofcontents
\newpage

%======================================================================
\section{The Logarithmic Vector Space}
\label{sec:log-vector}
%======================================================================

From the baseless logarithm framework (see companion document \texttt{baseless\_log.tex}), we have the logarithmic space $\mathcal{L}$ with:
\begin{itemize}
\item Elements: $\log x$ for $x \in \mathbb{R}^{>0}$
\item Addition: $(\log x) + (\log y) = \log(xy)$
\item Scalar multiplication: $c \cdot (\log x) = \log(x^c)$ for $c \in \mathbb{R}$
\end{itemize}

\begin{definition}[Logarithmic Vector]
A \textbf{logarithmic vector} is an element $\log x \in \mathcal{L}$. We write vectors in bold: $\mathbf{v} = \log x$.
\end{definition}

\begin{proposition}[$\mathcal{L}$ is a Real Vector Space]
The logarithmic space $(\mathcal{L}, +, \cdot)$ with the operations above is a real vector space isomorphic to $(\mathbb{R}, +, \cdot)$.
\end{proposition}

\begin{proof}
Define $\phi: \mathcal{L} \to \mathbb{R}$ by $\phi(\log x) = \ln x$. Then:
\begin{align*}
\phi((\log x) + (\log y)) &= \phi(\log(xy)) = \ln(xy) = \ln x + \ln y = \phi(\log x) + \phi(\log y) \\
\phi(c \cdot (\log x)) &= \phi(\log(x^c)) = \ln(x^c) = c \ln x = c \cdot \phi(\log x)
\end{align*}
So $\phi$ is a linear isomorphism.
\end{proof}

%======================================================================
\section{Bases and Coordinates}
\label{sec:bases}
%======================================================================

\begin{definition}[Logarithmic Basis]
A \textbf{basis} for $\mathcal{L}$ is a single nonzero element $\log b \in \mathcal{L}$ (since $\mathcal{L}$ is one-dimensional). The coordinate of $\log x$ with respect to $\log b$ is:
\[
v_b = \frac{\log x}{\log b} = \log_b x \in \mathbb{R}.
\]
\end{definition}

\begin{definition}[Coordinate Map]
The \textbf{coordinate map} $\phi_b: \mathcal{L} \to \mathbb{R}$ with respect to basis $\log b$ is:
\[
\phi_b(\log x) = \log_b x = \frac{\log x}{\log b}.
\]
\end{definition}

\begin{theorem}[Coordinate Independence]
For any two bases $\log b_1, \log b_2 \in \mathcal{L}$, the coordinate maps are related by:
\[
\phi_{b_1}(\log x) = \frac{\log b_2}{\log b_1} \cdot \phi_{b_2}(\log x).
\]
That is, $\phi_{b_1} = \frac{\log b_2}{\log b_1} \circ \phi_{b_2}$.
\end{theorem}

\begin{proof}
\[
\phi_{b_1}(\log x) = \frac{\log x}{\log b_1} = \frac{\log b_2}{\log b_1} \cdot \frac{\log x}{\log b_2} = \frac{\log b_2}{\log b_1} \cdot \phi_{b_2}(\log x). \qedhere
\]
\end{proof}

\begin{remark}
The factor $\frac{\log b_2}{\log b_1} = \log_{b_1}(b_2)$ is the \textbf{transition coefficient} between coordinate systems, analogous to the Jacobian $\frac{dx'}{dx}$ in differential geometry.
\end{remark}

%======================================================================
\section{Vectors as Logarithms of Translation Operators}
\label{sec:vectors-as-logs}
%======================================================================

We now reverse the correspondence: geometric vectors can be interpreted as logarithms.

\begin{definition}[Translation Operator]
In differential geometry, the \textbf{translation operator} $T^{\mathbf{v}}$ for a vector $\mathbf{v} = v_x \partial_x + v_y \partial_y$ is:
\[
T^{\mathbf{v}} = e^{\mathbf{v}} = e^{v_x \partial_x + v_y \partial_y}.
\]
This acts on functions as: $T^{\mathbf{v}} f(x,y) = f(x + v_x, y + v_y)$.
\end{definition}

\begin{proposition}[Vector as Logarithm of Translation]
The geometric vector $\mathbf{v}$ is the logarithm of its translation operator:
\[
\mathbf{v} = \ln T^{\mathbf{v}}.
\]
\end{proposition}

\begin{proof}
By definition, $T^{\mathbf{v}} = e^{\mathbf{v}}$, so $\ln T^{\mathbf{v}} = \mathbf{v}$.
\end{proof}

\begin{remark}
This means the space of geometric vectors is isomorphic to the space of translation operators via the exponential/logarithm map. The logarithm here is literally the inverse of the exponential map in Lie theory.
\end{remark}

\begin{theorem}[Decomposition in Coordinates]
For a vector $\mathbf{v} = v_x \partial_x + v_y \partial_y$, the translation operator decomposes as:
\[
T^{\mathbf{v}} = e^{v_x \partial_x} \cdot e^{v_y \partial_y} = T_x^{v_x} \cdot T_y^{v_y},
\]
where $T_x^{v_x} = e^{v_x \partial_x}$ is translation in the $x$-coordinate by $v_x$.
\end{theorem}

\begin{proof}
In flat space, the partial derivative operators commute: $[\partial_x, \partial_y] = 0$. Therefore:
\[
e^{v_x \partial_x + v_y \partial_y} = e^{v_x \partial_x} \cdot e^{v_y \partial_y}. \qedhere
\]
\end{proof}

\begin{remark}
Taking logarithms: $\mathbf{v} = \ln T^{\mathbf{v}} = v_x \ln T_x + v_y \ln T_y = v_x \partial_x + v_y \partial_y$. The components $v_x, v_y$ are the ``coordinates'' of $\mathbf{v}$ in the basis $\{\partial_x, \partial_y\}$.
\end{remark}

%======================================================================
\section{The Exponential Map and Logarithmic Structure}
\label{sec:exponential}
%======================================================================

\begin{definition}[Exponential Map]
The \textbf{exponential map} $\exp: \mathcal{L} \to \mathbb{R}^{>0}$ is defined by:
\[
\exp(\log x) = x.
\]
This is a group isomorphism from $(\mathcal{L}, +)$ to $(\mathbb{R}^{>0}, \times)$.
\end{definition}

\begin{proposition}[Exponential-Logarithm Duality]
The following diagram commutes:
\[
\begin{array}{ccc}
\mathcal{L} & \xrightarrow{\exp} & \mathbb{R}^{>0} \\
\downarrow_{c \cdot} & & \downarrow_{(\cdot)^c} \\
\mathcal{L} & \xrightarrow{\exp} & \mathbb{R}^{>0}
\end{array}
\]
That is, $\exp(c \cdot \log x) = (\exp(\log x))^c = x^c$.
\end{proposition}

\begin{proof}
\[
\exp(c \cdot \log x) = \exp(\log(x^c)) = x^c = (\exp(\log x))^c. \qedhere
\]
\end{proof}

%======================================================================
\section{Summary: The Complete Correspondence}
\label{sec:summary}
%======================================================================

\begin{table}[h]
\centering
\begin{tabular}{|l|l|}
\hline
\textbf{Logarithmic Structure} & \textbf{Vector Structure} \\
\hline
$\log x$ (baseless log) & $\mathbf{v}$ (geometric vector) \\
$\log b$ (choice of base/unit) & $\mathbf{e}$ (choice of basis) \\
$\log_b x = \frac{\log x}{\log b}$ (coordinate) & $v_e = \frac{\mathbf{v}}{\mathbf{e}}$ (component) \\
$\exp(\log x) = x$ (exponential) & $e^{\mathbf{v}} = T^{\mathbf{v}}$ (translation) \\
$\ln T^{\mathbf{v}} = \mathbf{v}$ (logarithm) & $\ln T^{\mathbf{v}} = \mathbf{v}$ (inverse exp) \\
Change of base $b_1 \to b_2$ & Change of coordinates $\mathbf{e}_1 \to \mathbf{e}_2$ \\
Conversion factor $\log_{b_1}(b_2)$ & Transition matrix $\frac{\mathbf{e}_2}{\mathbf{e}_1}$ \\
\hline
\end{tabular}
\caption{Complete correspondence between logarithmic and vector structures.}
\end{table}

\begin{theorem}[Main Isomorphism]
There is a natural isomorphism:
\[
\Phi: (\mathbb{R}^{>0}, \times) \xrightarrow{\sim} (\mathbb{R}, +) \xrightarrow{\sim} \text{Vect}(\mathbb{R}^2)
\]
where:
\begin{enumerate}
\item $\Phi(x) = \log x$ (baseless logarithm)
\item The middle isomorphism sends $t \mapsto t \partial_x$ (a one-dimensional vector)
\item Composition gives: $x \mapsto (\ln x) \partial_x$
\end{enumerate}
\end{theorem}

\end{document}
